\section{Research Agenda and Conclusion}
\label{sec:agenda}

The agenda shifts from a universal repackaging score to reviewable claims about artifact relations, production history, authority, and behavior. Detector evaluations should state the conclusion enabled and the premise supplied elsewhere.

\subsection{Direction-labeled, transformation-complete benchmarks}

Benchmarks should distinguish direct copy, common predecessor, shared library or template, legitimate version, independent generation, and authorized variant. They should record parent/child digests, ordered transformations, first-party/component regions, signing identities, chronology, and label provenance, while separating exact-history controlled cases from realistic field cases and retaining an unknown state.

Android pipelines combine shrinking, desugaring, generation, resource processing, cross-platform compilation, native code, and protection. Because order changes observability, lawful benchmarks should retain per-stage artifacts and report which relation survives each composition.

\subsection{Provenance-to-APK and component-to-behavior binding}

An Android provenance profile should terminate at the analyzed APK, binding source revision, resolved dependencies, generation/build configuration, builder, signer, publication event, and digest. Sensitive details may remain controlled, but commitments and predecessor relations must be checkable.

Declared composition should be reconciled with the binary. Reports must separate family, candidate version, partial inclusion, modification, reachability, and unknowns; reproducible builds, binary comparison, and transparency should act as independent checks rather than repeated declarations.

\subsection{Auditable generated transformations}

AI-mediated development needs artifact-level provenance. Each accepted patch or file set should bind a parent revision to a downstream product and record the model/service boundary, material retrieval/tool context, review decision, output digest, and unavailable evidence. Research must determine how to commit to sensitive context, retain sufficient agent history, cluster sibling generations without asserting copying, verify package recommendations, and represent model or tool drift when replay is impossible.

\subsection{Open-world, temporal, and uncertainty-aware evaluation}

Relatives, developers, libraries, markets, prompt families, and agent trajectories must not leak across random splits. Evaluations should hold out time and entity families, report candidate volume and abstention, keep unknown histories unknown, and distinguish absent, undetected, out-of-scope, and unavailable evidence.

Because keys, roles, packages, policies, and vulnerabilities change, decisions should be reproducible at a stated time and revisable through append-only records and explicit policy versions.

\subsection{Decision-centered assurance tooling}

Tooling should ingest artifact, component, provenance, signature, transparency, policy, and behavior evidence; check digest/predecessor bindings; and report supported, contradicted, and unresolved subclaims rather than one trust score.

Evaluation should follow retrieval, localization, adjudication, policy action, and appeal. With candidate population, policy, time, and review budget fixed, comparisons should test artifact-only, process-only, and bound cross-layer evidence and report precision/recall, abstention, reviewer effort, binding success, and localized disagreements. A decision log should also preserve the evidence available at the time, the policy version, the reason for abstention or escalation, and any later evidence that changed the outcome. Such revision tests whether the system supports accountable correction rather than only one-shot classification.

\subsection{Stewardship and practice}

AndroZoo and RePack show the value of shared data~\cite{androzoo,repack}. Resources should preserve label provenance, transformations, licenses, errors, and historical denominators; constrained evaluations can separate public synthetic, redistributable open-source, metadata-only field, and controlled-access strata.

Stores should separate retrieval, lineage, authorization, and behavior review. Developers should retain source-to-APK/dependency bindings, protect signing roles, preserve analyzed artifacts, and record accepted generated changes. Researchers should report the claim ceiling rather than overstate direction or authority.

\subsection{Conclusion}

Android lineage comprises identity, composition, resemblance, derivation, process, authority, and behavior. Repackaging, TPL, supply-chain, and AI-development research illuminate different links; none should be promoted beyond its observation.

A defensible assurance case binds source and dependencies to build steps, steps to an APK digest, the digest to an authorized transparent release, and the exact artifact to component and behavior analyses. Artifact comparison remains essential when records are absent or disputed; provenance supplies direction and accountability. This claim-centered structure is a stronger target than a universal repackaging score.
